%!TEX root = TheoretischePhysikC.tex

\section{Vorlesung: Harmonischer Oszillator}
Der Energie eines \emph{klassischen} eindimensionalen Oszillators ist gegeben durch
\begin{equation}
	E = \frac{1}{2m}p^2 + \frac{m\omega^2}{2}x^2,
\end{equation}
wobei $\omega$ die Frequenz des klassischen Oszillators mit dem gleichen Potential
ist.

Nach unserer Quantisierungstabelle lautet der Hamiltonoperator des quantumechanischen Oszillators:
\begin{equation}
	\widehat{H} = \frac{1}{2m}\widehat{p}^2 + \frac{m\omega^2}{2}\widehat{x}^2.
\end{equation}	


Der Oszillator tritt in vielen Kontexten der Physik auf: nämlich immer dann, wenn man kleine Bewegungen um eine Gleichgewichtslage näherungsweise untersucht. Ein freies Quantenfeld lässt sich als System von unendlich
vielen Oszillatoren analysieren. Es lohnt sich jedenfalls, dieses System genau kennenzulernen. Deshalb werden wir auch mehrere Betrachtungsweisen untersuchen.

Die direkteste Methode, das Spektrum zu erhalten, basiert auf der Beobachtung, dass sich $H$ (im wesentlichen) als Quadrat einer komplexen Linearkombination von $\widehat{p}$ und $\widehat{x}$ schreiben lässt. Wir setzen
\begin{equation}
	\widehat{a} = \sqrt{\frac{m\omega}{2\hbar}}\widehat{x} + i \sqrt{\frac{1}{2m\omega\hbar}}\widehat{p},
\end{equation}
sodass
\begin{align}
	\widehat{a} \widehat{a}^\dag &= \frac{m\omega}{2\hbar}\widehat{x}^2 + \frac{1}{2m\omega\hbar}\widehat{p}^2 + \frac{i}{2\hbar} (\widehat{p}\widehat{x} - \widehat{x}\widehat{p}) \\
	\widehat{a}^\dag \widehat{a} &= \frac{m\omega}{2\hbar}\widehat{x}^2 + \frac{1}{2m\omega\hbar}\widehat{p}^2 - \frac{i}{2\hbar} (\widehat{p}\widehat{x} - \widehat{x}\widehat{p}).
\end{align}
Unter Verwendung der kanonischen Vertauschungsrelationen $[\widehat{p},\widehat{x}] = -i\hbar\mathbb{I}$ ist damit
\begin{equation}\label{eq:ccra}
	\eqbox{ [\widehat{a}, \widehat{a}^\dag] = \mathbb{I}}
\end{equation}
und
\begin{equation}\label{eq:qho}
	\eqbox{\widehat{H} = \hbar \omega \left(\widehat{a}^\dag \widehat{a} + \tfrac{1}{2}\mathbb{I}\right).}
\end{equation}
Zunächst ist aus dieser Form von $H$ klar, dass $H \ge \hbar\omega/2 \mathbb{I}$. Der Oszillator ist
also niemals ,,in Ruhe``: $x=0$ und $p = 0$ können wegen der Unschärferelation nicht beide beliebig genau präpariert werden. 

Ist die ,,Nullpunktsenergie`` $\hbar \omega/2$ selbst ein Eigenwert? Der zugehörige Eigenvektor $|\phi_0\rangle$ müsste die Gleichung 
\begin{equation}
	\widehat{a}^\dag \widehat{a}|\phi_0\rangle = 0,
\end{equation}
also $\|\widehat{a}|\phi_0\rangle\|=0$, also $\widehat{a}|\phi_0\rangle = 0$ erfüllen. In der Ortsdarstellung ist
dies die gewöhnliche Differentialgleichung $\langle x|\widehat{a}|\phi_0\rangle = 0$:
\begin{equation}
	\sqrt{\frac{m\omega}{2\hbar}}x\phi_0(x) + i\sqrt{\frac{1}{2m\omega\hbar}} \frac{\hbar}{i}\frac{d}{dx} \phi_0(x) = 0,
\end{equation}
also 
\begin{equation}
	\frac{d\phi_0}{dx}  = -\frac{m\omega}{\hbar} \phi_0,
\end{equation}
mit der Lösung
\begin{equation}
	\phi_0(x) = Ne^{-\frac{m\omega}{\hbar} \frac{x^2}{2}},
\end{equation}
wobei $N$ die Normierung 
\begin{equation}
	N = \left(\frac{m\omega}{\pi\hbar}\right)^{\frac{1}{4}}
\end{equation}
ist.


In der Ableitung der Unschärferelation hatten wir genau wie hier die \emph{Positivität} das Operators $\widehat{a}^\dag \widehat{a}$ benutzt. Der Fall der Gleichheit führt auch dort auf die Bedingung $\widehat{a}|\Psi\rangle = 0$. Der Oszillator-Grundzustand ist also genau ein Zustand minimaler Unschärfe.

Für die weitere Rechnung brauchen wir die Relationen
\begin{align}
	[\widehat{H},\widehat{a}^\dag] &= \hbar \omega \widehat{a}^\dag \\
	[\widehat{H},\widehat{a}] &= -\hbar \omega \widehat{a}, 
\end{align}
die aus (\ref{eq:ccra}) und (\ref{eq:qho}) folgen. Wenn nun $\widehat{H}|\Psi\rangle = E|\Psi\rangle$ ein Eigenvektor ist, dann ist 
\begin{equation}
	\widehat{H}\widehat{a}^\dag|\Psi\rangle = (E  + \hbar \omega)\widehat{a}^\dag|\Psi\rangle
\end{equation}
Also erfüllt $\widehat{a}^\dag|\Psi\rangle$ die Eigenwertgleichung zum Eigenwert $E + \hbar \omega$. Wegen 
\begin{equation}
	\|\widehat{a}^\dag|\Psi\rangle\|^2 = \langle \Psi|\widehat{a}\widehat{a}^\dag|\Psi\rangle = \langle \Psi|\widehat{a}^\dag\widehat{a}+\mathbb{I}|\Psi\rangle \ge \langle \Psi|\Psi\rangle \ge 0
\end{equation}
kann $\widehat{a}^\dag|\Psi\rangle$ auch nicht Null sein: wir bekommen
also wirklich einen neuen Eigenvektor. Startend von $|\phi_0\rangle$ erhalten wir so die
Eigenvektoren 
\begin{equation}\label{eq:qhoeigenstates}
	\eqbox{|\phi_n\rangle = \frac{(\widehat{a}^\dag)^n|\phi_0\rangle}{\|(\widehat{a}^\dag)^n|\phi_0\rangle \|}}
\end{equation}
und Eigenwerte
\begin{equation}
	\eqbox{E_n = \hbar\omega(n+\tfrac{1}{2}).}
\end{equation}

Wegen dieses Verfahrens heißen übrigens die Operatoren $\widehat{a}$ und $\widehat{a}^\dag$ auch \emph{Leiteroperatoren}, bzw. ,,Absteigeoperator`` und ,,Aufsteigeoperator``. Den Absteigeoperator $\widehat{a}$ verwenden wir, um zu zeigen, dass es auch keine weiteren Eigenwerte gibt. Denn wenn $|\Psi\rangle$ ein Eigenvektor zum Eigenwert $E$ ist, erfüllt $(\widehat{a})^n|\Psi\rangle$ die Eigenwertgleichung zum Eigenwert $E − n\hbar \omega$. Für genügend große $n$ wird dieser Ausdruck negativ, kann also kein Eigenwert des positiven Operators
$H$ sein. Für große $n$ muss also $(\widehat{a})^n|\Psi\rangle=0$ sein. Sei $n$ die kleinste Zahl, für die dies gilt. Dann ist $|\Phi\rangle = (\widehat{a})^{n-1}|\Psi\rangle$ eine von Null verschiedene Lösung
der Gleichung $\widehat{a}|\Phi\rangle = 0$, also nach unserer obigen Rechnung proportional zu $|\phi_0\rangle$. Nach der Vertauschungsrelation für $\widehat{a}$ und $\widehat{a}^\dag$ ist $|\Psi\rangle$ dann ein Vielfaches von
$(\widehat{a}^\dag)^n|\Phi_0\rangle$. Die angegebene vollständige Liste von Eigenfunktionen ist also eine Basis des Hilbertraums.

Die Norm $\|(\widehat{a}^\dag)^n|\Phi_0\rangle\|$ lässt sich ebenfalls allein aus der Vertauschungsrelation
berechnen. Es ist
\begin{equation}
	\|\widehat{a}^\dag|\phi_n\rangle\|^2 = \langle\phi_n|\widehat{a}\widehat{a}^\dag|\phi_n\rangle =  \langle\phi_n|\left(\frac{1}{\hbar\omega}\widehat{H}+\frac12\mathbb{I}\right)|\phi_n\rangle=(n+1)\|\phi_n\|^2 = (n+1).
\end{equation}
Damit ist 
\begin{equation}
	\widehat{a}^\dag|\phi_n\rangle = \sqrt{n+1}|\phi_{n+1}\rangle,
\end{equation}
und (durch Induktion)
\begin{equation}
	|\phi_n\rangle = \frac{1}{\sqrt{n!}}(\widehat{a}^\dag)^n|\Phi_0\rangle.
\end{equation}

Da die $|\phi_n\rangle$ eine Basis des Hilbertraumes bilden, lassen sich alle Operatoren durch ihre Matrixelemente bezüglich dieser Basis ausdrücken. Das ist immer dann sinnvoll, wenn zeitabhängige Größen untersucht werden sollen, da ja dann die Zeitentwicklungsoperatoren Multiplikationsoperatoren sind. Wir bekommen
(mit der Matrixnotation $[\widehat{C}]_{nm} = \langle\phi_n|\widehat{C}|\phi_m\rangle$)
\begin{equation}
	[\widehat{a}^\dag]_{nm} = \sqrt{n}\delta_{n,m+1} \quad \text{und} \quad [\widehat{a}]_{nm} = \sqrt{m}\delta_{n+1,m}.
\end{equation}
Mit 
\begin{equation}
	\widehat{p} = \frac{\sqrt{2m\omega\hbar}}{2i}(\widehat{a}-\widehat{a}^\dag) 
\end{equation}
bekommt man daraus unmittelbar die Matrix für $\widehat{p}$ und entsprechend die Matrix für $\widehat{x}$. Hingemalt sind es:
\begin{equation}
	\widehat{p} = \frac{\sqrt{2m\omega\hbar}}{2i}
	\begin{pmatrix}
		0 & \sqrt{1} & 0 & 0 & 0 & 0 & \hdots \\
		-\sqrt{1} & 0 & \sqrt{2} & 0 & 0 & 0 & \\
		0 & -\sqrt{2} & 0 & \sqrt{3} & 0 & 0 & \\
		0 & 0 & -\sqrt{3} & 0 & \sqrt{4} & 0 & \\
		0 & 0 & 0 & -\sqrt{4} & 0 & \sqrt{5} & \\
		0 & 0 & 0 & 0 & -\sqrt{5} & 0 & \\
		\vdots & & & & & & \ddots \\
	\end{pmatrix}
\end{equation}
und
\begin{equation}
	\widehat{x} = \sqrt{\frac{\hbar}{2m\omega}}
	\begin{pmatrix}
		0 & \sqrt{1} & 0 & 0 & 0 & 0 & \hdots \\
		\sqrt{1} & 0 & \sqrt{2} & 0 & 0 & 0 & \\
		0 & \sqrt{2} & 0 & \sqrt{3} & 0 & 0 & \\
		0 & 0 & \sqrt{3} & 0 & \sqrt{4} & 0 & \\
		0 & 0 & 0 & \sqrt{4} & 0 & \sqrt{5} & \\
		0 & 0 & 0 & 0 & \sqrt{5} & 0 & \\
		\vdots & & & & & & \ddots \\
	\end{pmatrix}.
\end{equation}

Ein wichtiger Vorteil der oben erwähnten \emph{algebraischen Methode} der Leiteroperatoren ist,
dass die Ergebnisse sofort in jeder Darstellung des Hilbertraums zur Verfügung stehen. Zum Beispiel liefert die Formel (\ref{eq:qhoeigenstates}) die Wellenfunktionen/Eigenfunktionen $\phi_n(x) \in L^2(\mathbb{R})$ in der Ortsdarstellung, sobald $\widehat{a}$ konkret in dieser Darstellung eingesetzt wird. Durch Induktion sieht man, dass 
\begin{equation}
	\phi_n(x) = c_n \mathrm{H}_n(\lambda x)  e^{-\frac{\lambda^2 x^2}{2}}
\end{equation}
mit
\begin{equation}
	\lambda = \sqrt{\frac{m\omega}{\hbar}}\quad \text{und} \quad c_n = \sqrt{\frac{\lambda}{2^n n!\sqrt{\pi}}},
\end{equation}
wobei
$\mathrm{H}_n(\lambda x)$ die sogenannten \emph{Hermite-Polynome} sind. Die ersten $3$ Hermite-Polynome
lauten
\begin{align}
	\mathrm{H}_0(z) &= 1,\\
	\mathrm{H}_1(z) &= 2z,\\
	\mathrm{H}_2(z) &= 4z^2-2,
\end{align}
usw.

Mit der Diagonalisierung des Hamiltonoperators haben wir natürlich auch detaillierte Informationen über die Zeitentwicklung des Oszillators gewonnen. Wie der klassische Oszillator ist auch der quantenmechanische periodisch mit der Periode $T = 2\pi/\omega$. Das folgt allein schon aus der Liste der Eigenwerte: die Eigenwerte des Zeitentwicklungsoperators $U_T = e^{-\frac{iT}{\hbar} \widehat{H}}$ sind ja nach Funktionalkalkül
\begin{equation}
	e^{-\frac{iT}{\hbar} \hbar\omega(n+\frac12)} = e^{-2\pi i (n+\frac12)} = -1,
\end{equation}
also unabhängig von $n$, und $U_T = −\mathbb{I}$ ist ein Vielfaches des $\mathbb{I}$-Operators.

Damit gilt für die Zeitabhängigkeit aller Erwartungswerte, nach Postulat 5 also für die Funktionen der Form $f(t) = \tr(U_t \rho U_t^\dag A)$:
\begin{equation}
	f(t+T) = \tr(U_tU_T \rho U_T^\dag U_t^\dag A) =f(t).
\end{equation}	
Die Analogie mit der klassischen Entwicklung geht aber noch viel weiter. Am einfachsten sieht man das hier, indem man die Zeitabhängigkeit des Leiteroperators $\widehat{a}$ berechnet. Es gilt 
\begin{equation}
	\frac{d}{dt} U_t^\dag \widehat{a} U_t = U_t^\dag \left(\frac{i}{\hbar} [\widehat{H},\widehat{a}] \right) U_t = -i\omega U_t^\dag \widehat{a} U_t
\end{equation}
also
\begin{equation}
	U_t^\dag \widehat{a} U_t = e^{-i\omega t}\widehat{a}.
\end{equation}
Dabei findet man, dass die quantenmechanischen Orts-und Impulsoperatoren genau den klassischen Bewegungsgleichungen folgen:
\begin{equation}
	\begin{split}
		\widehat{p}(t) &= U_t^\dag \widehat{p} U_t  = \cos(\omega t) \widehat{p}  - m\omega\sin(\omega t) \widehat{x} \\
		\widehat{x}(t) &= U_t^\dag \widehat{x} U_t  = \frac{1}{m\omega}\sin(\omega t) \widehat{p}  + \cos(\omega t) \widehat{x}
	\end{split}
\end{equation}